\documentclass[UTF8,fontset=macnew,11pt]{ctexart}

\usepackage[a4paper,margin=2.25cm]{geometry}
\usepackage{amsmath,amssymb}
\usepackage{booktabs}
\usepackage{tabularx}
\usepackage{array}
\usepackage{xcolor}
\usepackage{hyperref}

\hypersetup{
  colorlinks=true,
  linkcolor=blue!50!black,
  citecolor=blue!50!black,
  urlcolor=blue!50!black
}

\setlength{\parskip}{0.35em}
\setlength{\parindent}{2em}
\setlength{\emergencystretch}{2em}

\newcommand{\method}{Resource-NMCTS}
\newcommand{\pareto}{Pareto-Resource-NMCTS}
\newcommand{\score}{\ensuremath{\mathrm{score}}}
\newcommand{\anf}{\ensuremath{\mathrm{ANF}}}
\newcommand{\fprm}{\ensuremath{\mathrm{FPRM}}}
\newcommand{\mcts}{\ensuremath{\mathrm{MCTS}}}

\title{面向资源约束量子布尔函数综合的结构级神经搜索方法}
\author{中文技术论文稿 v5}
\date{2026年7月8日}

\begin{document}
\pagestyle{plain}
\maketitle

\begin{abstract}
量子布尔函数 oracle 综合需要把经典谓词嵌入可逆量子线路，其逻辑层 T-count、CNOT、深度和辅助比特数会直接影响容错资源估计。直接代数正规形（algebraic normal form, \anf{}）综合稳定但容易重复计算高次数 AND 结构；SSHR 等 CNOT-oriented 方法提供了重要小规模参照，却不等同于本文研究的资源加权搜索问题。本文将 oracle 综合建模为 \anf{}/固定极性 Reed--Muller（fixed-polarity Reed--Muller, \fprm{}）项集合上的组合优化，提出由资源感知分解、布尔环线性因子筛选、神经动作先验、蒙特卡洛树搜索、结构级 depth policy 与 screen-gate 组成的逻辑层综合框架。实验显示，在 $n\leq6$ 的 177 个传统函数上，\pareto{} 相对 ESOP cube beam 获得 174/0/3 的 score 胜/负/平，平均 score 降低 36.09\%；相对 weighted ESOP-MILP 获得 167/3/7，平均 score 降低 29.84\%。在 $n=18$ 的 12 个随机 \anf{} 函数上，adaptive depth-2 布尔环筛选相对单层筛选获得 11/0/1，平均 score 降低 6.47\%；完整 \method{} 相对 adaptive screen 进一步降低 23.85\%，说明中等高维仍需要深搜索。在 $n=20$ 的 6 个边界函数上，adaptive depth-2 screen 相对单层筛选降低 7.47\% 的平均 score，并使 \method{} 相对 AND-direct ANF 降低 11.80\%；screen-gated \method{} 在保持全部资源持平的同时，将平均运行时间从 77.10 s 降至 31.18 s。进一步的结构级 depth policy 在 held-out $n=20$ 项集上相对单层和 depth-1 screen 分别降低 6.57\% 和 2.57\% 的平均 score，并比全 depth adaptive 评估节省 34.71\% 时间；保守 depth-2 skip guard 则在另一组 held-out $n=20$ 项集上与固定 depth-2 screen 全部持平，并相对全 depth adaptive 节省 2.87\% 时间。上述结果说明 AI 模块已经从动作排序推进到结构选择，但固定 depth-2 screen 仍是很强的确定性基线，因此本文把当前贡献界定为资源搜索框架与结构级门控证据，而非高维神经先验已经独立主导质量提升。
\end{abstract}

\noindent\textbf{关键词：}量子 oracle 综合；布尔函数综合；神经蒙特卡洛树搜索；结构级策略；布尔环；ANF；FPRM

\section{引言}

给定布尔函数 $f:\{0,1\}^n\rightarrow\{0,1\}$，量子 oracle 通常实现为
\begin{equation}
  |x\rangle |y\rangle \mapsto |x\rangle |y\oplus f(x)\rangle .
  \label{eq:oracle}
\end{equation}
这类 oracle 是 Grover 搜索、量子密码分析和若干量子算法子程序中的基础组件。由于容错量子计算中非 Clifford 门代价高，oracle 的 T-count、CNOT、逻辑深度和辅助比特数不能被视为后端细节，而是资源估计能否可信的核心组成。

\anf{} 是最直接的 oracle 输入表示。若
\begin{equation}
  f(x)=\bigoplus_{m\in\mathcal{M}}\prod_{i\in m}x_i ,
  \label{eq:anf}
\end{equation}
则每个单项式都可被翻译为可逆乘法结构，并异或到输出位。该路线语义透明，但会忽略不同单项式之间的公共 AND 结构、线性奇偶结构和极性重写机会。对于随机高维项集合，重复计算会迅速推高 T-count。

已有 XAG、ROS、BDD 和 SSHR 等方法从不同中间表示处理 oracle synthesis 问题\cite{meuli2022xag,meuli2020ros,wille2009bdd,zheng2025sshr}。本文不从 SSHR 的 parallelotope cover 空间出发，也不把主目标设为 CNOT-only 最优；SSHR 在本文中是 CNOT-oriented small-scale baseline。本文的核心问题是：能否把量子布尔函数综合写成一个资源加权搜索问题，并让 AI 模块在组合搜索中承担明确、可验证的作用。

本文的一句话论证是：在逻辑层量子布尔函数 oracle 综合中，\anf{}/\fprm{} 项集合上的资源感知搜索，配合布尔环线性因子筛选和结构级神经门控，可以在传统函数和 $n=18,20$ 高维随机函数上降低 T-count 与加权逻辑资源；当前边界是高维 learned prior 的独立质量收益仍不足，结论也不覆盖硬件 mapping、routing 或物理噪声。

\section{问题定义与资源模型}

本文输入为布尔函数的真值表或 \anf{} 表示，输出为实现式~\eqref{eq:oracle} 的逻辑层可逆线路。每条候选线路都必须通过 truth-table 级语义验证。候选选择使用统一逻辑资源分数
\begin{equation}
  \score = T + 0.04\,\mathrm{CNOT} + 0.015\,\mathrm{depth}
         + 0.01\,\mathrm{gates} + 2\,\mathrm{ancilla}.
  \label{eq:score}
\end{equation}
式~\eqref{eq:score} 是可复现实验中的排序准则，不代表具体硬件平台的真实物理成本。因此本文同时报告 T-count、CNOT、depth、gates 与 peak ancilla，避免单一 score 掩盖资源 trade-off。

\method{} 的状态是当前尚未综合的项集合 $\mathcal{M}$。一次动作选择一个可计算、可反算的局部因子结构，将项集合分解为 quotient 子问题与 rest 子问题，再递归生成线路。候选动作包括公共单项式因子、固定极性重写后的公共项、CNOT-only 线性因子、布尔环线性因子，以及由神经模型排序或扩展的 root action。相对于直接在真值表上训练黑盒策略，该状态保留了布尔代数结构，并使每一步都可翻译为线路片段。

\section{方法}

\subsection{资源感知搜索主干}

直接 \anf{} 综合逐项计算单项式。若多个单项式共享高次数因子，则可以先计算该因子，再用较低成本的 quotient 子问题复用它。\fprm{} 极性搜索进一步改变项集合形状，使原始 \anf{} 中不明显的公共结构变得显式。\mcts{} 在固定预算内探索候选组合，神经动作先验根据候选覆盖率、项数变化、平均次数、因子大小、即时资源收益和子问题规模等特征对动作排序。

由于 learned prior 在高维随机项集合上存在负迁移风险，当前实现采用 baseline-preserving guard。也就是说，确定性候选和学习候选同时产生，最终只接受不劣于 baseline 的结果。这个设计降低了 AI 模块导致资源退化的风险，也使实验结果更适合写成保守、可复现的综合框架。

\subsection{布尔环线性因子筛选}

旧 linear-pair 分支主要搜索形如
\begin{equation}
  (x_i\oplus x_j\oplus\cdots)\cdot h(x)
  \label{eq:linear_factor}
\end{equation}
的局部结构，并要求线性因子变量与 quotient $h(x)$ 变量不相交。这个约束便于实现，但会漏掉 square-free Boolean ring 中合法的重叠变量结构。布尔环线性因子允许二者共享变量，并在展开时使用 $x_i^2=x_i$ 与 GF(2) 偶数项消去。例如
\begin{equation}
  (x_i\oplus x_j)x_i m = x_i m \oplus x_i x_j m .
  \label{eq:boolean_ring}
\end{equation}
线路层面仍可先用 CNOT 计算线性奇偶量，再用该辅助量控制 quotient 子线路。因此，该扩展不是单纯扩大 beam 宽度，而是把原先被实现约束排除的代数候选重新纳入搜索。

\subsection{结构级 depth policy}

高维 $n=18,20$ 实验显示，固定单层 screen、depth-1 screen 和 depth-2 screen 的收益随项集合结构变化。为把 AI 贡献从动作排序推进到结构选择，本文训练轻量策略网络预测应使用 single、depth-1 还是 depth-2 screen。输入特征包括项数、次数分布、变量覆盖密度、二元共现密度，以及 root Boolean-ring action 的数量和质量统计。监督标签不是人工规则，而是实际运行三种 screen 后按式~\eqref{eq:score} 选择的 oracle depth。进一步地，本文训练一个保守 depth-2 skip guard，专门判断 depth-1 是否足以与固定 depth-2 screen 持平；阈值在训练和验证切片上按 zero-false-skip 约束选择。

这个模块的目标不是替代强 baseline，而是减少全 depth adaptive 评估成本，并学习 ``何时需要更深筛选''。因此本文把 depth policy 写成结构级 AI 证据，而不是最终主质量贡献。

\subsection{Screen-gate 运行时门控}

在 $n=20$ 边界切片中，完整 \method{} 的最终资源与 adaptive screen 持平，但运行时间更长。为控制这一长尾，本文训练一个可解释 decision stump，判断 adaptive screen 是否已经足够、是否可以跳过昂贵的 Resource tail。当前模型由 $n=18,20$ 的 18 个配对样本训练得到，学到的规则等价于在 $n\geq19$ 时优先跳过 Resource tail。由于训练样本很少，这一模块只被解释为运行时门控证据，不被解释为普适质量模型。

\section{实验设置}

实验分为四组。第一组使用 $n\leq6$ 的 177 个传统函数，与 direct ANF、AND-direct ANF、ESOP cube beam、ESOP-MILP、SSHR-H、ABC 和 BDD baseline 比较。第二组使用 $n=18$ 随机 \anf{} 函数，观察 adaptive screen 与完整 \method{} 的差距。第三组使用 $n=20$ 边界函数，检验在 FPRM root beam 和 fast linear-pair 超时区域中，depth-2 screen 是否仍能给出可运行改进。第四组使用 $n=14,16,18$ 合成项集训练 depth policy 与 depth-2 skip guard，并在 held-out $n=20$ 项集上测试结构选择能力。

所有 paired comparison 只纳入函数名匹配、语义正确且未 timeout 的结果。出现 timeout 的方法保留为运行边界证据，但不参与正确结果的均值比较。

\section{结果}

\subsection{传统小规模函数}

表~\ref{tab:traditional} 给出 $n\leq6$ 传统函数上的平均资源。相对 direct ANF，\pareto{} 平均 T-count 降低 72.25\%，平均 score 降低 67.80\%。相对 ESOP cube beam，\pareto{} 获得 174/0/3 的 score 胜/负/平，平均 score 降低 36.09\%。相对 weighted ESOP-MILP，获得 167/3/7，平均 score 降低 29.84\%。

\begin{table}[t]
\centering
\small
\caption{$n\leq6$ 传统函数切片上的平均逻辑资源。}
\label{tab:traditional}
\begin{tabular}{lrrrrr}
\toprule
方法 & 平均 T & 平均 CNOT & 平均 depth & 平均 ancilla & 平均 score \\
\midrule
Direct ANF & 184.1 & 134.2 & 134.6 & 1.33 & 194.3 \\
AND-direct ANF & 101.0 & 166.5 & 166.9 & 2.18 & 115.2 \\
\method{} & 43.9 & 89.9 & 94.5 & 1.92 & 53.2 \\
\pareto{} & \textbf{40.4} & 83.0 & \textbf{88.0} & 2.03 & \textbf{49.6} \\
ESOP-MILP & 83.6 & 133.5 & 159.6 & 2.32 & 96.7 \\
SSHR-H & 81.0 & \textbf{67.6} & 88.7 & 1.36 & 88.2 \\
\bottomrule
\end{tabular}
\end{table}

该表同时暴露了边界：SSHR-H 的平均 CNOT 最低，因此本文不能主张 CNOT-only 全面支配。更准确的主张是，本文方法在当前逻辑层加权资源模型下用一定 CNOT 与辅助比特 trade-off 换取更低 T-count 和更低 score。

\subsection{\texorpdfstring{$n=18$}{n=18} 随机 ANF}

表~\ref{tab:n18} 显示，adaptive depth-2 Boolean screen 相对单层 Boolean screen 获得 11/0/1 的 score 胜/负/平，平均 score 从 11349.06 降至 10670.94，相对降低 6.47\%。同时，完整 \method{} 平均 score 为 7315.62，相对 adaptive screen 进一步降低 23.85\%。这说明 $n=18$ 上快速 screen 是有效候选生成器，但不能替代完整搜索。

\begin{table}[t]
\centering
\small
\caption{$n=18$ 随机 \anf{} 函数上的 adaptive screen 与完整 \method{}。}
\label{tab:n18}
\begin{tabular}{lrrrrr}
\toprule
方法 & 平均 T & 平均 CNOT & 平均 depth & 平均 score & 平均时间(s) \\
\midrule
Single Boolean screen & 10389.33 & 16268.42 & 16268.50 & 11349.06 & 0.82 \\
Adaptive depth-2 screen & 9767.00 & 15301.58 & 15301.67 & 10670.94 & 4.39 \\
\method{} & \textbf{6639.33} & \textbf{11380.92} & \textbf{11382.17} & \textbf{7315.62} & 226.18 \\
\bottomrule
\end{tabular}
\end{table}

\subsection{\texorpdfstring{$n=20$}{n=20} 边界切片}

$n=20$ 切片更接近当前实现的扩展边界。此前 FPRM root beam 和 fast linear-pair 在该切片中均出现 300 s task timeout。表~\ref{tab:n20} 显示，adaptive depth-2 screen 平均 score 为 21331.24，低于单层 screen 的 22695.15 和 depth-1 screen 的 21988.52。相对单层 screen，adaptive depth-2 screen 获得 5/0/1，平均 score 降低 7.47\%。集成该分支后，\method{} 相对 AND-direct ANF 获得 5/0/1，平均 score 降低 11.80\%。

\begin{table}[t]
\centering
\scriptsize
\caption{$n=20$ 边界函数上的 screen 与 \method{} 结果。}
\label{tab:n20}
\begin{tabularx}{\linewidth}{>{\raggedright\arraybackslash}Xrrrrr}
\toprule
方法 & T & CNOT & depth & score & 时间(s) \\
\midrule
Direct ANF & 42944.00 & 19600.33 & 19600.33 & 44037.63 & 10.36 \\
AND-direct ANF & 21516.67 & 33635.67 & 33635.67 & 23491.15 & 10.41 \\
Single screen & 20786.00 & 32492.50 & 32492.50 & 22695.15 & 10.70 \\
Depth-1 screen & 20138.00 & 31480.50 & 31480.50 & 21988.52 & 12.03 \\
Adaptive depth-2 screen & \textbf{19535.33} & \textbf{30542.50} & \textbf{30542.50} & \textbf{21331.24} & 20.67 \\
\method{} & \textbf{19535.33} & \textbf{30542.50} & \textbf{30542.50} & \textbf{21331.24} & 77.10 \\
Screen-gated \method{} & \textbf{19535.33} & \textbf{30542.50} & \textbf{30542.50} & \textbf{21331.24} & 31.18 \\
\bottomrule
\end{tabularx}
\end{table}

Screen-gated \method{} 与完整 \method{} 的 T-count、CNOT、depth、peak ancilla 和 score 全部持平，但平均运行时间降低 61.31\%。这不是新的质量收益，而是高维长尾控制收益。该结果支持把门控机制写进方法框架，但不能据此断言 Resource tail 在所有高维函数上都可跳过，因为 $n=18$ 中完整搜索仍明显强于 screen。

\subsection{结构级 depth policy}

表~\ref{tab:policy} 给出 held-out $n=20$ 项集上的 depth policy 结果。策略网络在 $n=14,16,18$ 上训练，在 48 个 $n=20$ 测试项集上评估。它相对固定 single screen 获得 42/0/6，平均 score 降低 6.57\%；相对 depth-1 screen 获得 34/0/14，平均 score 降低 2.57\%。相对全 depth adaptive 评估，policy 的平均 score 只高 0.28\%，但平均运行时间降低 34.71\%。

\begin{table}[t]
\centering
\small
\caption{结构级 depth policy 在 held-out $n=20$ 项集上的比较。}
\label{tab:policy}
\begin{tabular}{lrrrr}
\toprule
比较 & 函数数 & score 胜/负/平 & 平均 $\Delta$ score & 平均 $\Delta$ time \\
\midrule
Policy vs single screen & 48 & 42/0/6 & $-6.57\%$ & $+2224.00\%$ \\
Policy vs depth-1 screen & 48 & 34/0/14 & $-2.57\%$ & $+257.27\%$ \\
Policy vs depth-2 screen & 48 & 0/5/43 & $+0.28\%$ & $-10.85\%$ \\
Policy vs all-depth adaptive & 48 & 0/5/43 & $+0.28\%$ & $-34.71\%$ \\
\bottomrule
\end{tabular}
\end{table}

这组结果的正向意义在于，AI 模块不再只是对 root action 排序，而是开始学习结构选择。它的限制也很清楚：固定 depth-2 screen 在 score 上仍略优。为去掉这一质量缺口，本文进一步训练了保守 depth-2 skip guard。表~\ref{tab:depth_guard} 显示，该 guard 在 train、validation 和另一组 held-out $n=20$ test 中均没有 false skip；测试集上 48 个样本全部与固定 depth-2 screen score 持平，相对全 depth adaptive 评估节省 2.87\% 时间。

\begin{table}[t]
\centering
\small
\caption{保守 depth-2 skip guard 的 zero-loss 结果。时间变化以全 depth adaptive 为 baseline。}
\label{tab:depth_guard}
\begin{tabular}{lrrrrr}
\toprule
切片 & 函数数 & skip depth-2 & false skip & score 胜/负/平 & 平均 $\Delta$ time \\
\midrule
Train & 240 & 22 & 0 & 0/0/240 & $-4.48\%$ \\
Validation & 72 & 3 & 0 & 0/0/72 & $-1.61\%$ \\
Held-out test & 48 & 3 & 0 & 0/0/48 & $-2.87\%$ \\
\bottomrule
\end{tabular}
\end{table}

这个 guard 的意义是把结构级 policy 相对固定 depth-2 的 score 缺口从 $+0.28\%$ 降为 $0$，但其收益仍然小。由于它需要先评估浅层 screen，再决定是否回退到 depth-2，测试集上相对单独固定 depth-2 screen 反而慢 39.40\%。因此它应被解释为 ``质量安全门控方向已打通''，而不是最终运行时优化已经完成。

\section{讨论}

本文最稳妥的贡献表述有三点。第一，方法主线独立于 SSHR。本文不是把 SSHR 改成 AI 版本，而是在 \anf{}/\fprm{} 项集合上做资源加权搜索；SSHR 只作为 CNOT-oriented 小规模 baseline。第二，布尔环线性因子提供了明确的结构扩展。它允许线性因子与 quotient 共享变量，捕获旧实现约束排除的候选，并在 $n=18,20$ 上稳定改善单层 screen。第三，结构级 AI 已有初步正向证据。Depth policy 学会了何时需要更深 screen，保守 depth-2 guard 消除了相对固定 depth-2 的 score 缺口，screen-gate 也能在 $n=20$ 边界切片上减少完整 Resource tail 的运行时间。

但当前稿件仍不能声称高维神经搜索已经充分成功。原因是：$n=18$ 上完整 \method{} 明显强于 adaptive screen，说明中等高维仍需要更深搜索；$n=20$ 上完整 \method{} 与 adaptive screen 资源持平，说明该切片的收益主要来自结构筛选而非 MCTS tail；保守 depth-2 guard 虽然保持了 fixed depth-2 的 score，但跳过样本只有 3/48，运行时收益仍小。审稿表达上，应把这些结果写成 ``结构级 AI 证据与下一步突破方向''，而不是 ``最终高维质量模型''。

\section{下一步明显提升目标}

若以明显提升作为阶段结束点，下一步应集中在三项可量化目标上。第一，扩大 guard-aware depth policy 的跳过覆盖率：目标是在 held-out $n=20$ 上继续保持不劣于固定 depth-2 screen 的 score，同时把相对 all-depth adaptive 的时间优势从 2.87\% 提高到至少 15\%。第二，把结构级策略扩展为分支选择器：不仅选择 screen 深度，还预测是否进入 \fprm{} polarity search、完整 \method{} 或 screen-gated tail。第三，补充函数级线路案例和外部 reversible-toolchain 对比，使论文从表格结果走向可解释的综合机制。

\section{结论}

本文提出一种面向资源约束量子布尔函数综合的结构级神经搜索方法。该方法以 \anf{}/\fprm{} 项集合为状态，以逻辑资源 score 为选择准则，用布尔环线性因子、神经先验、\mcts{}、baseline guard、depth policy 和 screen-gate 共同生成候选线路。实验表明，\pareto{} 在 $n\leq6$ 传统函数上显著降低 T-count 与 score；adaptive depth-2 screen 在 $n=18,20$ 高维随机函数上稳定改善单层 screen；screen-gate 在 $n=20$ 边界函数上保持资源不变并显著减少运行时间；保守 depth-2 guard 初步实现了不劣于固定 depth-2 的结构级选择。当前结果已经形成一条独立于 SSHR 的论文主线，但要达到更强投稿标准，还需要把结构级 AI 从小幅安全门控推进到更大覆盖率的质量收益。

\begin{thebibliography}{99}

\bibitem{meuli2019multiplicative}
G.~Meuli, M.~Soeken, E.~Campbell, M.~Roetteler, and G.~De Micheli.
\newblock The role of multiplicative complexity in compiling low T-count oracle circuits.
\newblock In \emph{Proceedings of ICCAD}, 2019.

\bibitem{meuli2022xag}
G.~Meuli, M.~Soeken, and G.~De Micheli.
\newblock Xor-And-Inverter Graphs for quantum compilation.
\newblock \emph{npj Quantum Information}, 8:7, 2022.

\bibitem{meuli2020ros}
G.~Meuli, M.~Soeken, M.~Roetteler, and G.~De Micheli.
\newblock ROS: Resource-constrained oracle synthesis for quantum computers.
\newblock \emph{Electronic Proceedings in Theoretical Computer Science}, 318:119--130, 2020.

\bibitem{wille2009bdd}
R.~Wille and R.~Drechsler.
\newblock BDD-based synthesis of reversible logic for large functions.
\newblock In \emph{Proceedings of DAC}, pp.~270--275, 2009.

\bibitem{zheng2025sshr}
Q.~Zheng, Y.~Xu, J.~Zhang, Z.~Su, and S.~Zheng.
\newblock CNOT oriented synthesis for small-scale Boolean functions using spatial structures of parallelotopes.
\newblock In \emph{Proceedings of ICCAD}, 2025.

\bibitem{tsaras2024shortcircuit}
D.~Tsaras, A.~Grosnit, L.~Chen, Z.~Xie, H.~Bou-Ammar, and M.~Yuan.
\newblock ShortCircuit: AlphaZero-driven synthesis of Boolean circuits.
\newblock \emph{arXiv:2408.09858}, 2024.

\bibitem{rietsch2024unitary}
S.~Rietsch, A.~Y. Dubey, C.~Ufrecht, M.~Periyasamy, A.~Plinge, C.~Mutschler, and D.~D. Scherer.
\newblock Unitary synthesis of Clifford+T circuits with reinforcement learning.
\newblock In \emph{IEEE International Conference on Quantum Computing and Engineering}, pp.~824--835, 2024.

\bibitem{riu2025rlzx}
J.~Riu, J.~Nogu{\'e}, G.~Vilaplana, A.~Garcia-Saez, and M.~P. Estarellas.
\newblock Reinforcement learning based quantum circuit optimization via ZX-calculus.
\newblock \emph{Quantum}, 9:1758, 2025.

\bibitem{machiya2026monteq}
M.~Machiya, M.~Menickelly, P.~Hovland, and J.~Liu.
\newblock MonteQ: A Monte Carlo tree search based quantum circuit synthesis framework.
\newblock \emph{arXiv:2604.19029}, 2026.

\end{thebibliography}

\end{document}
